\documentclass[12pt,letterpaper]{article}

\usepackage[margin=1in]{geometry}
\usepackage{amsmath,amssymb}
\usepackage{fancyhdr}

\setlength{\parindent}{0pt}
\setlength{\parskip}{0.8em}
\setlength{\headheight}{15pt}

\pagestyle{fancy}
\fancyhf{}
\fancyhead[L]{FIT3155 SAMPLE EXAM PAPER}
\fancyfoot[C]{Page \thepage{} of 22}
\renewcommand{\headrulewidth}{0pt}
\renewcommand{\footrulewidth}{0pt}

\newcommand{\marksbox}[1]{%
  \vfill
  \begin{flushright}
    \fbox{\parbox[c][1.1cm][c]{1.3cm}{\centering #1}}
  \end{flushright}
}

\newcommand{\questiontitle}[2]{%
  \textbf{Question #1}\hfill\textbf{(#2 marks)}
}

\newcommand{\workingpage}{%
  \newpage
  \begin{center}
    (Page intentionally left blank for working space)
  \end{center}
}

\begin{document}

\begin{titlepage}
  \thispagestyle{empty}
  \vspace*{0.18\textheight}
  {\Large FIT3155 \hfill SAMPLE EXAM PAPER\par}
  \vspace{2.2cm}
  \begin{center}
    {\LARGE\bfseries FIT3155 SAMPLE EXAM PAPER\par}
    \vspace{0.8cm}
    {\large Generated Sample Exam 04\par}
  \end{center}
\end{titlepage}

\setcounter{page}{2}

\section*{INSTRUCTIONS}

\begin{itemize}
  \item This exam is worth 60 marks.
\end{itemize}

\begin{center}
  \textbf{General exam technique}
\end{center}

\begin{itemize}
  \item Some students lose marks by not attempting all questions.

  \item Suppose you are likely to get 3/6 on a question after spending say 10 minutes on it.
  Spending another 10 minutes on the same question may gain you at most 3 additional
  marks. In contrast, spending those 10 minutes on a new question could earn you another
  6 marks. The key point is: do not get stuck on one question.

  \item A good strategy is to first answer the questions you are most confident about before
  attempting the others. In other words, answering questions strictly serially during an
  exam can be suboptimal.

  \item Do not waste time answering an unrelated question in the hope of gaining partial marks.
  Marking will be based on the question that was actually asked.

  \item Where necessary, justify your answers with clear and concise explanations, without being
  overly wordy.
\end{itemize}

\newpage

\questiontitle{1}{6=2+4}

For any string $S[1 \ldots n]$, define the \emph{Z-suffix value}
$Zsuf_i(S)$, for each position $1 \leq i < n$, as the length of the longest
substring of $S$ \emph{ending} at position $i$ that is identical to a suffix
of $S$.

Example: for $S[1 \ldots 7] = \texttt{aabxaab}$, $Zsuf_3(S) = 3$, since
$S[1 \ldots 3] = \texttt{aab}$ is identical to the suffix
$S[5 \ldots 7] = \texttt{aab}$.

\begin{enumerate}
  \item[(i)] Compute $Zsuf_i(S)$ for every position $1 \leq i < n$ of the
  string
  \[
    S[1 \ldots 6] = \texttt{babcab}.
  \]

  \item[(ii)] Write pseudocode describing an efficient algorithm that takes
  any input string $S[1 \ldots n]$ and outputs all of its Z-suffix values.
  State the worst-case time complexity of your algorithm.

  (For this problem, you may assume that \texttt{zalg} is available as a
  built-in function that takes a string as input and returns the array of its
  Z-values. Do not write pseudocode for \texttt{zalg}.)
\end{enumerate}

\marksbox{6}

\workingpage
\newpage

\questiontitle{2}{6}

In the Boyer--Moore algorithm for exact matching of the pattern
$pat[1 \ldots m]$ against the text $txt[1 \ldots n]$, suppose that during some
iteration the pattern is aligned with $txt[j \ldots j+m-1]$, and the
right-to-left scan finds that \emph{every} character matches, i.e., an
occurrence of $pat$ has been found starting at position $j$ of the text.

To continue searching for further occurrences, the pattern is then shifted
$m - mp[2]$ positions to the right, where $mp[2]$ denotes the length of the
longest suffix of $pat[2 \ldots m]$ that is identical to a prefix of $pat$.

Prove that this shift is safe, i.e., such a shift can never skip $pat$ past
any valid occurrence in $txt$.

\marksbox{6}

\workingpage
\newpage

\questiontitle{3}{6=2+4}

Let $S[1 \ldots n]$ be a string whose final character $S[n] = \$$ is a unique
terminal symbol, lexicographically smaller than every other character. Let
$SA[1 \ldots n]$ be the suffix array of $S$, let $L[1 \ldots n]$ be the
Burrows--Wheeler Transform of $S$, and let $F[1 \ldots n]$ denote the first
column of the sorted cyclic rotation matrix (i.e., the characters of $S$ in
sorted order).

\begin{enumerate}
  \item[(i)] Write pseudocode that computes $L[1 \ldots n]$ given only $S$ and
  $SA$ as input. State the worst-case time complexity of your procedure.

  \item[(ii)] Prove the following property: for every character $c$ and every
  $i$, the $i$-th occurrence of $c$ in $L$ (reading top to bottom) and the
  $i$-th occurrence of $c$ in $F$ (reading top to bottom) correspond to the
  same occurrence of $c$ in $S$.

  (This is the rank-preservation property underpinning the LF mapping used in
  BWT inversion and backward search.)
\end{enumerate}

\marksbox{6}

\workingpage
\newpage

\questiontitle{4}{6=3+3}

Assume that a suffix tree for $S[1 \ldots n]\$$ has \emph{already} been
constructed, with each edge label stored as a pair of indices
$(start, end)$ into the string. You may assume the alphabet is of constant
size.

\begin{enumerate}
  \item[(i)] Describe a preprocessing step that takes $O(n)$ time, after
  which the \emph{number} of occurrences in $S$ of any query pattern
  $pat[1 \ldots m]$ can be reported in $O(m)$ time. Justify why the count
  your method reports is correct, and why the stated time bounds hold.

  \item[(ii)] Describe an $O(n)$-time algorithm that finds the longest
  \emph{repeated} substring of $S$, i.e., the longest substring that occurs
  at least twice in $S$. Justify the correctness of your algorithm.
\end{enumerate}

\marksbox{6}

\workingpage
\newpage

\questiontitle{5}{6=3+3}

\begin{enumerate}
  \item[(i)] A text of length 100 contains exactly five distinct symbols,
  with the following frequencies:
  \[
  \begin{array}{c|ccccc}
    \text{symbol} & \texttt{a} & \texttt{b} & \texttt{c} & \texttt{d} & \texttt{e} \\
    \hline
    \text{frequency} & 38 & 25 & 20 & 10 & 7
  \end{array}
  \]
  Construct a Huffman code for these symbols, showing the order in which
  subtrees are merged. Give a valid codeword for each symbol, and compute the
  total length (in bits) of the resulting Huffman encoding of the text.

  \item[(ii)] Work out the Elias-Omega binary codeword for the positive
  integer $100$. Show each length component of your codeword.
\end{enumerate}

\marksbox{6}

\workingpage
\newpage

\questiontitle{6}{6=3+3}

\begin{enumerate}
  \item[(i)] Write pseudocode for an efficient algorithm that computes
  \[
    a^{b} \bmod n
  \]
  for positive integers $a$, $b$, $n$, using the method of repeated squaring.
  State, as a function of $b$, the number of modular multiplications your
  algorithm performs in the worst case.

  \item[(ii)] Using repeated squaring, compute
  \[
    5^{21} \bmod 23.
  \]
  Show the value of each successive squaring modulo $23$, and indicate which
  of those values are multiplied together to obtain the final answer.
\end{enumerate}

\marksbox{6}

\workingpage
\newpage

\questiontitle{7}{6=3+3}

Let $x$ be any node of degree $k$ in a Fibonacci heap, and let
$y_1, y_2, \ldots, y_k$ denote the current children of $x$, in the order in
which they were linked to $x$ (earliest first). Recall that nodes become
children of other nodes only through the consolidation step (which links two
roots of \emph{equal} degree), and that a non-root node is cut from its parent
as soon as it loses a second child.

\begin{enumerate}
  \item[(i)] Prove that $\deg(y_i) \geq i - 2$ for all $2 \leq i \leq k$.

  \item[(ii)] Using part (i), prove that the subtree rooted at any node of
  degree $k$ contains at least $F_{k+2}$ nodes, where $F_0 = 0$, $F_1 = 1$,
  and $F_{j} = F_{j-1} + F_{j-2}$. Hence conclude that the maximum degree of
  any node in an $n$-node Fibonacci heap is $O(\log n)$.

  (You may assume without proof that
  $\sum_{i=0}^{k} F_i = F_{k+2} - 1$ and that $F_{k+2} \geq \phi^{k}$, where
  $\phi = (1 + \sqrt{5})/2$.)
\end{enumerate}

\marksbox{6}

\workingpage
\newpage

\questiontitle{8}{6}

A stack data structure supports the following three operations:

\begin{itemize}
  \item \textsc{Push}$(x)$: pushes $x$ onto the stack, with actual cost $1$;
  \item \textsc{Pop}$()$: pops the top element, with actual cost $1$;
  \item \textsc{MultiPop}$(k)$: pops $\min(k, \mathit{size})$ elements off the
  stack, with actual cost equal to the number of elements actually popped.
\end{itemize}

The stack is initially empty.

Using the \emph{accounting} (banker's) method of amortized analysis, determine
the amortized cost of each of the three operations. State the charge you
assign to each operation, state the credit invariant that your scheme
maintains, and verify that the total stored credit can never become negative
over any sequence of operations starting from the empty stack.

\marksbox{6}

\workingpage
\newpage

\questiontitle{9}{6=3+3}

Consider the following B-tree with minimum degree $t = 2$, so that every node
stores between $t - 1 = 1$ and $2t - 1 = 3$ keys (the root is exempt from the
lower bound):

\begin{itemize}
  \item the root contains the keys $[20, 40]$;
  \item its three children are, from left to right, $[10, 15]$, $[25, 30]$,
  and $[50]$.
\end{itemize}

Use the following deletion conventions. A key in an internal node is deleted
by replacing it with its in-order \emph{predecessor} (the largest key in the
child subtree immediately to its left), which is then deleted from its
original node. If a deletion would leave a non-root node with fewer than
$t - 1$ keys, then: (a) if an immediate sibling has more than $t - 1$ keys,
\emph{borrow} through the parent (the separating key in the parent moves down
into the deficient node, and the sibling's key adjacent to the separator
moves up into the parent); (b) otherwise, \emph{merge} the deficient node
with an immediate sibling and the separating key from the parent.

\begin{enumerate}
  \item[(i)] Delete key $50$. Draw the resulting B-tree, clearly indicating
  the keys in every node.

  \item[(ii)] From the B-tree obtained in part (i), delete key $20$. Draw the
  resulting B-tree, clearly indicating the keys in every node.
\end{enumerate}

\marksbox{6}

\workingpage
\newpage

\questiontitle{10}{6=3+3}

\begin{enumerate}
  \item[(i)] Convert the following linear program into the standard form used
  by the simplex method (a maximization objective, $\leq$ constraints only,
  and non-negative decision variables), and then rewrite each constraint as
  an equality by introducing slack variables. Do \emph{not} solve the
  program.
  \[
  \begin{array}{rl}
  \text{Minimize} & z = 2x_1 - 3x_2 \\
  \text{Subject to} & x_1 + x_2 = 10, \\
                    & x_1 - 2x_2 \geq -4, \\
                    & x_1 \geq 0,\quad x_2 \geq 0.
  \end{array}
  \]

  \item[(ii)] Consider the linear program:
  \[
  \begin{array}{rl}
  \text{Maximize} & z = 5x + 4y \\
  \text{Subject to} & 2x + y \leq 8, \\
                    & x + 3y \leq 9, \\
                    & x \geq 0,\quad y \geq 0.
  \end{array}
  \]
  Introduce slack variables and write down the \emph{initial} simplex
  tableau. State the basic feasible solution that this tableau represents.
  Then identify the \emph{entering} variable and, using the minimum-ratio
  test, the \emph{leaving} variable for the first pivot of the simplex
  method. Do \emph{not} perform the pivot.
\end{enumerate}

\marksbox{6}

\workingpage

\vfill
\begin{center}
  End of Exam
\end{center}

\end{document}
